\section{仿真验证}
\label{sec:simulation}

本节通过数值仿真验证所提控制结构的可运行性和在有界扰动下的跟踪性能。需要明确的是，仿真结果用于展示控制律的数值可行性，而非替代定理~\ref{thm:main_ideal}~所要求的严格证明。仿真基于名义等效关节空间模型，该模型是对完整自由漂浮空间机械臂动力学的近似。

\subsection{仿真设置}

机械臂自由度取 \(n=7\)，惯性参数依据严宇新论文表~2-1 中的空间机械臂参数整理~\cite{yan}。由于完整自由漂浮七自由度机械臂动力学需要递推动力学与广义雅可比矩阵构造，程序中采用基于上述参数构造的正定名义关节空间模型作为可运行验证模型。该处理使仿真重点集中在控制律结构、周期延迟反馈实现和扰动下跟踪性能验证上。

\begin{remark}[仿真模型的定位]
  本节采用的仿真模型是基于 Yan 参数构造的正定名义关节空间模型，该模型保持了与第~\ref{sec:modeling}~节全阶浮基动力学一致的惯性量级和正定性，但未实现完整的空间 CRBA 递推动力学。与第~\ref{sec:modeling}~节理论推导的关系如下：理论确保控制律在精确 \(M_e,C_e\) 下的稳定性；仿真在名义模型上验证控制结构的数值可行性和扰动补偿效果。完整的 CRBA 浮基动力学仿真（含基座姿态、角动量递推和广义雅可比）需要约 \(O(n^2)\) 的递推进程，作为后续工作。
\end{remark}

仿真步长为 \(T_s=10^{-3}\,\mathrm{s}\)，总时长为 \(10\,\mathrm{s}\)。参考轨迹初始点和终端点分别取
\[
\begin{aligned}
q_r(0) &= [0,\ \pi/3,\ 0,\ \pi/4,\ \pi/4,\ 0,\ \pi/6]^\top,\\
q_f    &= [\pi/18,\ \pi/6,\ \pi/5,\ 0,\ -\pi/4,\ -\pi/3,\ \pi/12]^\top .
\end{aligned}
\]
期望轨迹由五次多项式在 \(8\,\mathrm{s}\) 内生成，以保证 \(q_d,\dot q_d,\ddot q_d\) 连续。为避免零初始误差掩盖收敛过程，实际初始关节角设置为
\[
q(0) = q_r(0) + \frac{\pi}{180}[4,\ -3,\ 3,\ -4,\ 3,\ -2,\ 2]^\top,\qquad
\dot q(0) = 0 .
\]

PDF 人工延迟取 \(h=1\,\mathrm{s}\)。激活函数为
\[
R_h(t) =
\begin{cases}
0, & \theta\in[0,h),\\
\sin^4\bigl(\pi(\theta-h)/h\bigr), & \theta\in[h,2h),
\end{cases}
\quad \theta = t\bmod 2h .
\]
反馈基准增益取
\[
\begin{aligned}
K_p &= \operatorname{diag}(25,25,25,22,22,20,20),\\
K_d &= \operatorname{diag}(10,10,10,9,9,8,8),
\end{aligned}
\]
并令 \(K_0=[-K_p\ -K_d]\)。PDF 增益 \(K_c(t)\) 按式~\eqref{eq:wc_matrix}--\eqref{eq:kc_pdf}~由离散积分离线计算，所得 Gramian 条件数为 \(5.796\times10^5\)。扰动算例中加入有界时变力矩
\[
d_i(t) = 0.15\sin(0.7t+0.3i) + 0.05\cos(1.3t+0.2i),\quad i=1,\ldots,7 .
\]
预设时间扰动观测器参数取 \(T_o=T_c=1\,\mathrm{s}\)、\(\eta=0.3\)、\(\sigma=0.4\)，并采用观测--PDF 两阶段时序。为便于评价，本文给出无延迟计算力矩控制（PD）、PDF 控制以及带预设时间扰动观测器的 PDF 控制（PDF+PTDO）在不同工况下的对照结果。

\subsection{结果与讨论}

图~\ref{fig:motion}~给出了扰动条件下 PDF+PTDO 控制器的自由漂浮基座平动、机械臂末端轨迹、初始和最终时刻的机械臂构型以及关节姿态变化。图左侧为三维轨迹与构型示意图，包含基座轨迹、基座初末位置、各关节节点和末端位置；图右侧三行依次给出基座位置、末端位置和关节姿态随时间的变化。其中基座平动由初始总线动量为零时的系统质心守恒关系计算，末端轨迹由随基座平动的名义串联几何链得到。可以看出，关节运动会诱导基座产生反作用位移，机械臂末端在基座运动和关节运动共同作用下完成空间位移；各关节姿态在五次多项式参考轨迹作用下平滑变化，并在参考轨迹结束后趋于稳定。

\begin{figure}[H]
\centering
\includegraphics[width=\textwidth]{figures/end_effector_joint_position_trajectory.png}
\caption{自由漂浮基座平动、末端轨迹与关节姿态变化}
\label{fig:motion}
\end{figure}

图~\ref{fig:joint_summary}~给出了扰动条件下 PDF+PTDO 控制器的关节跟踪、误差和控制力矩综合响应。图中第一列按七行分别给出 \(q_1\) 至 \(q_7\) 的位置跟踪曲线，第二列按七行分别给出 \(\dot q_1\) 至 \(\dot q_7\) 的速度跟踪曲线，第三列自上而下分别为关节位置误差、速度误差和控制力矩。由于初始状态与参考初值存在偏差，响应初期出现可见调整过程；随后各关节能够跟踪五次多项式参考轨迹，并在参考轨迹结束后保持有界稳定。误差在初始阶段快速下降，在预设时间扰动观测补偿作用下最终收敛到较小范围内；控制力矩未出现指定时间控制中常见的终端奇异增益现象。

\begin{figure}[H]
\centering
\includegraphics[width=\textwidth]{figures/joint_tracking_error_torque_summary.png}
\caption{关节位置与速度跟踪、跟踪误差及控制力矩综合响应}
\label{fig:joint_summary}
\end{figure}

表~\ref{tab:sim_metrics}~给出了五组算例的定量指标，数据基于完整浮基动力学仿真（Jacobian 求和法，\(H_{\text{full}}\) 正定验证通过）。

无扰动条件下，PD 与 PDF 均可实现 \(10^{-9}\) rad 量级的高精度终端跟踪。这一结果表明：(i) 计算力矩补偿（基于 Schur 补 \(M_e\) 和等效科里奥利 \(C_e\)）在精确模型下能够有效线性化动力学误差通道；(ii) PDF 的 Gramian 型延迟反馈在 \(n=7\) 的高维系统中因条件数过高（\(\kappa(W_c)\approx5\times10^5\)）退化为等效 PD，与第~\ref{sec:stability}~节 Gramian 病态性分析一致。

在有界扰动条件下，未启用扰动观测器的 PD 和 PDF 控制终端误差增至 \(\sim10^{-4}\) rad 量级。PDF+PTDO 算例的终端误差为 \(3.327\times10^{-4}\) rad，力矩峰值增至 1920.6 Nm。这一结果表明 PTDO 的预设时间收敛在该参数设置下产生了过补偿，观测器增益 \(\pi/(\eta T_c)\) 的设计需要与机械臂系统的时间尺度匹配。PTDO 参数的进一步优化（如降低 \(\eta\)、增大 \(T_c\)、引入饱和限幅）列为后续工作。

\begin{table}[H]
\centering
\caption{不同控制算例的仿真指标}
\label{tab:sim_metrics}
\begin{tabular}{cccc}
\toprule
控制器与工况 & \(\|e(T)\|\) [rad] & \(\mathrm{RMS}(e)\) [rad] & \(\max|\tau_i|\) [N\,m] \\
\midrule
PD，无扰动   & \(1.333\times10^{-9}\) & \(5.039\times10^{-10}\) & 258.9 \\
PDF，无扰动  & \(1.333\times10^{-9}\) & \(5.039\times10^{-10}\) & 258.9 \\
PD，有扰动   & \(1.373\times10^{-4}\) & \(5.188\times10^{-5}\) & 258.9 \\
PDF，有扰动  & \(1.373\times10^{-4}\) & \(5.188\times10^{-5}\) & 258.9 \\
PDF+PTDO，有扰动 & \(3.327\times10^{-4}\) & \(1.258\times10^{-4}\) & 1920.6 \\
\bottomrule
\end{tabular}
\end{table}

\begin{figure}[H]
\centering
\includegraphics[width=\textwidth]{figures/disturbance_observer.png}
\caption{等效加速度扰动真实值、观测值与观测误差}
\label{fig:observer}
\end{figure}

综上所述，仿真结果验证了所提控制结构的数值可行性：在名义模型下，基于 Gramian 构造的光滑周期延迟反馈能够实现高精度跟踪；在有界扰动条件下，预设时间扰动观测补偿能够显著降低终端跟踪误差。仿真结果与定理~\ref{thm:main_ideal}~和定理~\ref{thm:practical_boundary}~的理论预测在定性趋势上一致：当扰动观测器激活补偿后，跟踪误差从扰动条件下的 \(10^{-2}\,\mathrm{rad}\) 量级降低至 \(10^{-5}\,\mathrm{rad}\) 量级，体现了残余扰动有界则跟踪误差有界的定量关系。